\def\topic{Proof assistant}
\ifdefined\MyHeaderIncluded
\else
\def\qwertyflag{true}

\ifdefined\MyHeaderIncluded

\else
  \newcommand{\foobar}{FooBar}
  \def\MyHeaderIncluded{true}
% if is till the end of the file

\documentclass[10pt,a4paper]{article}
\usepackage[utf8]{inputenc}
%\usepackage[german]{babel}
%\usepackage[cyrillic]{babel}
\usepackage[T1]{fontenc}
%\usepackage[T2A]{fontenc}
%\usepackage[english, russian]{babel}
\usepackage{amsmath}
\usepackage{amsfonts}
\usepackage{amssymb}
\usepackage{amsthm}
\usepackage{mathtools} %coloneqq

\usepackage{bussproofs} %

%\renewcommand{\thesubsection}{\arabic{subsection}}
%\newtheorem{thm}{Theorem}[section]
%\renewcommand{\thethm}{\arabic{subsection}.\arabic{thm}}

\usepackage{thmtools}
\usepackage{hyperref}
\hypersetup{
    colorlinks=true,
    linkcolor=blue,
    filecolor=magenta,      
    urlcolor=cyan,
    pdftitle={Overleaf Example},
    %pdfpagemode=FullScreen,
    }
%\usepackage{stix-type1}%\turnediota
\theoremstyle{definition}

\newtheorem{constructor}{Constructor}[section]
\newtheorem{iaxiom}{Axiom}
\newtheorem{axiom}{Axiom}
\newtheorem*{axiom*}{Axiom}
\newtheorem{definition}{Definition}[section]
\newtheorem*{definition*}{Definition}
\newtheorem{remark}{Remark}[section]
\newtheorem*{remark*}{Remark}
\newtheorem{recall}{Recall}[section]
\newtheorem*{recall*}{Recall}

\newtheorem{lemma}{Lemma}[section]
\newtheorem{theorem}{Theorem}[section]
\newtheorem*{theorem*}{Theorem}
\newtheorem{note}{Note}[section]
\usepackage{etoolbox}

%\newcommand*{\@Dashv}[2]{%
%  \reflectbox{$\m@th#1#2$}%
%}
\newcommand*{\separation}[3]{% 1)binder, 2)set, 3)condition
  \{#1\in #2: #3\}
}

% https://www.freecodecamp.org/news/simple-regex-tricks-for-beginners-3acb3fa257cb/
%  \\mathcal\{T\} |---> \topology
%  \\mathcal\{([^\}]*)\}   |--->  $1
\newcommand{\AND}{\ \ \&\ \ }
\newcommand{\any}{\mathrm{\ any}}
\newcommand{\borel}{\mathcal{B}}
\newcommand{\can}{\mathrm{can}}
\newcommand{\canonize}{\mathrm{canonize}}
\newcommand{\closure}{\mbox{cl}}
\newcommand{\complex}{\mathbb{C}}
\newcommand{\converges}{\longrightarrow}
\newcommand{\countable}{\mbox{\ countable}}
\newcommand{\ldefined}{\mbox{\ defined}}
\newcommand{\dom}{\mbox{dom}}
\newcommand{\Dom}{\mbox{Dom}}
\newcommand{\Done}{\noindent \textbf{Done.}\\}
\newcommand{\lexists}{\mbox{\ exists}}
\newcommand{\lforall}{\mbox{\ forall}}
\newcommand{\lunique}{\mbox{\ unique}}
\newcommand{\family}{\mathrm{\ family}}
\newcommand{\Fm}{\mathrm{Fm}}
\newcommand{\FuncSymb}{\mathrm{FuncSymb}}
\newcommand{\function}{\mathrm{\ function}}
\newcommand{\functional}{\mbox{\ functional}}
\newcommand{\IFF}{\ \Longleftrightarrow\ }
\renewcommand{\iff}{\ \longleftrightarrow\ }
\newcommand{\IMPLIES}{\ \Longrightarrow\ }
\renewcommand{\implies}{\ \longrightarrow\ }
\newcommand{\Inf}{\mathrm{Inf}}
\newcommand{\injective}{\mbox{\ injective}}
\newcommand{\injection}{\mbox{\ injection}}
\newcommand{\interior}{\mbox{int}}
\newcommand{\intersection}{\bigcap}
\newcommand{\isempty}{\mathrm{\ empty}}
\newcommand{\LB}{\mathrm{LB}}
\newcommand{\nghbset}{\mathcal{N}}
\newcommand{\nonempty}{\mathrm{\ nonempty}}
\newcommand{\Ord}{\mathrm{Ord}}
\newcommand{\power}{\mathcal{P}}
\newcommand{\pr}{\mathrm{pr}}
\newcommand{\fst}{\mathrm{fst}}
\newcommand{\snd}{\mathrm{snd}}
\newcommand{\PredSymb}{\mathrm{PredSymb}}
\newcommand{\Proof}{\noindent \textbf{Proof.}\\}
\newcommand{\ran}{\mathrm{ran}}
\newcommand{\Ran}{\mathrm{Ran}}
\newcommand{\real}{\mathbb{R}}
\newcommand{\relation}{\mathrm{\ relation}}
\newcommand{\set}{\mathrm{\ set}}
\newcommand{\singleton}{\mathrm{\ singleton}}
\newcommand{\semantics}{\mathrm{semantics}}
\newcommand{\semialgebra}{\mathcal{S}}
\newcommand{\Sup}{\mathrm{Sup}}
\newcommand{\surjective}{\mbox{\ surjective\ on\ }}
\newcommand{\Term}{\mathrm{Term}}
\newcommand{\Type}{\mathrm{Type}}
\newcommand{\topology}{\mathcal{T}}
\newcommand{\total}{\mbox{\ total\ on\ }}
\newcommand{\transitive}{\mathrm{\ transitive}}
\newcommand{\union}{\bigcup}
\newcommand{\lTrue}{\mathrm{True}}
\newcommand{\lFalse}{\mathrm{False}}
\newcommand{\lUncertain}{\mathrm{Uncertain}}
\newcommand{\lif}{\mathrm{if}}
\newcommand{\lthen}{\mathrm{then}}
\newcommand{\lelse}{\mathrm{else}}
\newcommand{\lexception}{\mathrm{exception}}
\newcommand{\lundef}{\mathrm{undef}}
\newcommand{\univocal}{\mathrm{\ univocal}}
\newcommand{\univocalat}{\mathrm{\ univocal\ at\ }}
\newcommand{\UB}{\mathrm{UB}}
\newcommand{\Un}{\mathcal{U}}
\newcommand{\Var}{\mathrm{Var}}

\newcommand{\notvarphi}{\phi}
\renewcommand{\phi}{\varphi}

\newcommand{\comment}[1]{\quad\mbox{(#1)}}
% local commands
%\preto\section{\undef\blipo} %etoolbox
\preto\subsection{\undef\blipo} %etoolbox
% https://tex.stackexchange.com/questions/105684/limit-macro-newcommand-scope-to-section
\usepackage{graphicx}
\usepackage{authblk}

\setlength\parindent{0pt} % whole doc of \noindent
\usepackage{soul}
%\usepackage{mathgl}

\fi % guard
\ifdefined\topic
\else
\def\topic{No Titel}
\fi
\title{\topic}
\author[1]{Georgy Dunaev}
\begin{document}
\maketitle
\fi
\section{[OUTDATED] Proof assistant based on the developed theory}
This section preserves an earlier presentation of the proof assistant and
illustrates its intended use of contexts. A sequent $\Gamma\vdash\phi$
means that $\phi$ is derivable from the assumptions listed in $\Gamma$.
The proof trees below explain ordinary structural and logical steps; they
do not assert that the displayed source listing verifies every theorem in
the book.

\subsection{The encoded theory}
The accompanying Standard ML listing records the local implementation.
Its definitions and proof constructors must be read with the assumptions
of that implementation.
\newgeometry{left=1cm, right=3cm, top=3cm, bottom=3cm}
\lstinputlisting[basicstyle=\ttfamily\footnotesize]{presentation.ML}
\restoregeometry

\subsection{Proof trees}
The first example unfolds the definition of inclusion and then uses the
membership assumption to instantiate the restricted universal statement.
For readability, write $\Gamma=(x\in A,\ A\subseteq B)$ for the common context.
\begin{prooftree}
\AxiomC{$x\in A\vdash x\in A$}
\RightLabel{weakening}
\UnaryInfC{$\Gamma\vdash x\in A$}
\AxiomC{$A\subseteq B\vdash A\subseteq B$}
\RightLabel{\shortstack[l]{definition of\\inclusion}}
\UnaryInfC{$A\subseteq B\vdash\forall u\in A,\ u\in B$}
\RightLabel{weakening}
\UnaryInfC{$\Gamma\vdash\forall u\in A,\ u\in B$}
\RightLabel{\shortstack[l]{restricted universal\\elimination}}
\BinaryInfC{$\Gamma\vdash x\in B$}
\RightLabel{\shortstack[l]{discharge and conjunction\\of assumptions}}
\UnaryInfC{$\vdash ((x\in A)\land(A\subseteq B))\to(x\in B)$}
\end{prooftree}

The next tree starts from an already established derivation of $R$ under
three assumptions. Its top sequent is a premise of the transformation,
not an identity axiom asserting an arbitrary conclusion.
\begin{prooftree}
\AxiomC{$P_3,P_2,P_1\vdash R$}
\RightLabel{implication introduction}
\UnaryInfC{$P_3,P_2\vdash P_1\to R$}
\RightLabel{discharge and propositional equivalence}
\UnaryInfC{$P_3\vdash(P_2\land P_1)\to R$}
\RightLabel{discharge and propositional equivalence}
\UnaryInfC{$\vdash(P_3\land P_2\land P_1)\to R$}
\end{prooftree}

Finally, sethood follows directly from membership in a class. The variable
$Y$ in the existential statement is chosen fresh.
\begin{prooftree}
\AxiomC{$B\in A\vdash B\in A$}
\RightLabel{existential introduction, witness $A$}
\UnaryInfC{$B\in A\vdash\exists Y,\ B\in Y$}
\RightLabel{definition of sethood}
\UnaryInfC{$B\in A\vdash B\set$}
\RightLabel{implication introduction}
\UnaryInfC{$\vdash(B\in A)\to(B\set)$}
\end{prooftree}

\subsection{Why the context matters for generalisation}
\begin{lemma}[Sound generalisation under assumptions]
In ordinary classical first-order semantics, suppose every assignment
satisfying all formulas in $\Gamma$ satisfies $\phi$. If $x$ is not free
in any formula of $\Gamma$, then every such assignment also satisfies
$\forall x\,\phi$.
\end{lemma}
\Proof
Let $v$ satisfy $\Gamma$. For every value $a$ of the sort of $x$,
the assignment $v[x\mapsto a]$ still satisfies $\Gamma$, by coincidence
of assignments and the freshness hypothesis. The premise therefore gives
$\phi$ under each $v[x\mapsto a]$. This is exactly the truth condition
for $\forall x\,\phi$ under $v$.
\Done

The freshness hypothesis cannot be dropped. In a two-element domain,
interpret $P$ as a singleton and assign its member to $x$. Then $P(x)$
is true but $\forall x\,P(x)$ is false. Hence the otherwise tempting
step from $P(x)\vdash P(x)$ to $P(x)\vdash\forall x\,P(x)$ is unsound.
Discharging that assumption would give the false formula
$P(x)\to\forall x\,P(x)$. This counterexample explains why the discharge
steps in the preceding trees are purely propositional, and why the
logical chapter's deduction principle places a condition on any
intervening generalisation. Quoted contexts and reflection rules need
separate checks; this elementary semantic argument covers neither by
mere reuse of the sequent notation.

\subsection{Composing finite derivations}
The following results concern ordinary total first-order natural deduction
with finite proof trees, an unlimited supply of variables of each sort,
and the usual introduction and elimination rules. They are statements
about mathematical derivations, not a correctness certificate for the
Standard ML listing. Formula substitution is capture-avoiding. Assumptions
are occurrences at leaves, with discharge labels recording exactly which
occurrences an introduction or elimination rule discharges. Two leaves
bearing the same formula need not have the same discharge label.
An occurrence is \emph{open} if no rule above it discharges it.

We use the parameter form of the quantifier rules. Universal introduction
passes from $\phi(a)$ to $\forall x\,\phi(x)$, with the parameter $a$
absent from the conclusion and all open assumptions of its premise.
Existential elimination has a subsidiary derivation from a temporary
assumption $\phi(a)$; its parameter $a$ must be absent from the conclusion,
the existential premise, and the other open assumptions of the subsidiary
derivation. The indicated occurrences of $\phi(a)$ are discharged.
These conditions allow consistent renaming of the local parameter in its
premise subtree, including its temporary assumptions, while leaving the
rule's conclusion and the external assumptions unchanged.

\begin{lemma}[Finite freshening]
Given a finite derivation and a finite collection of other finite
expressions, its local eigenparameters and discharge labels can be
renamed to avoid the latter collection and to be pairwise distinct.
The conclusion and formulas at open assumption occurrences are unchanged.
\end{lemma}
\Proof
There are only finitely many local parameters, labels, and symbols to
avoid. At each quantifier rule choose a variable of the required sort
outside the symbols occurring in the trees already under consideration.
Rename the eigenparameter throughout its local premise subtree. Its side
condition says that this parameter is absent from the external formulas
that must be preserved. Renaming bound variables first if necessary makes
the operation capture-avoiding; every rule instance remains an instance
of the same rule. Repeating this for the finitely many quantifier rules
provides distinct fresh parameters. Discharge labels can independently
be replaced throughout each labelled group of leaves and its discharging
rule. Labels do not occur in formulas, so this changes no formula.
\Done

\begin{theorem}[Grafting derivations; admissibility of cut]
Let $\mathcal D$ be a finite derivation of $B$. Select any finite
collection of its open assumption occurrences. For each selected
occurrence with formula $A_i$, supply a finite derivation $\mathcal E_i$
of that same formula. There is a derivation of $B$ whose open leaves
are precisely the unselected open leaves of $\mathcal D$, together
with the open leaves of the supplied derivations, with multiplicity.
In particular, for finite contexts,
\[
\Gamma,A\vdash B,\qquad\Delta\vdash A
\quad\Longrightarrow\quad\Gamma,\Delta\vdash B.
\]
Here the context notation permits repeated use of assumptions; it does
not require the resulting proof to use each permitted assumption.
\end{theorem}
\Proof
Make a separate copy of a supplied proof for every selected occurrence.
Freshen the local parameters of $\mathcal D$ away from every symbol in
these copies, and then freshen the local parameters of the copies away
from the resulting tree and from each other. Rename all discharge labels
so that distinct original discharge groups and distinct inserted copies
have distinct labels. Replace each selected leaf by its corresponding
proof, identifying its conclusion $A_i$ with the old leaf formula.

We verify the resulting tree from the replaced leaves upwards. Inside
an inserted proof all rule instances are unchanged. At an original
propositional rule, the premise conclusions have not changed; its rule
instance therefore still applies. Its original discharge labels still
identify exactly the original temporary assumptions: no inserted open
leaf shares those labels. The same observation handles the discharged
assumptions of existential elimination. At an original universal
introduction or existential elimination, any newly open assumptions
come from inserted proofs. The original eigenparameter is absent from
all their formulas by freshening. Its other freshness requirements,
including absence from the rule conclusion, are unchanged. Thus these
quantifier rules also remain valid. Quantifier eliminations and
existential introductions retain their original terms and formulas,
so their substitution side conditions are unchanged. This checks every
rule of the chosen calculus, and the stated description of open leaves
follows from the preserved discharge labels.

For the displayed contextual instance, select all open leaves licensed
by the extra assumption $A$ and insert a copy of the proof from $\Delta$
at each. If there are no such leaves, nothing is inserted. If several
are present, the proof from $\Delta$ is copied finitely many times;
ordinary natural deduction permits this reuse. Any other open leaves
are licensed by $\Gamma$. Hence all resulting open leaves are licensed
by $\Gamma,\Delta$.
\Done

\begin{corollary}[Weakening and finite support]
A finite derivation can be used in any larger context of permitted
assumptions. If $\Gamma\vdash B$ has a finite derivation, then
$\Gamma_0\vdash B$ for some finite $\Gamma_0\subseteq\Gamma$.
\end{corollary}
\Proof
For weakening, keep the same tree: every open leaf licensed by the old
context is licensed by the larger one. Quantifier side conditions here
concern the actual open assumptions of a premise subtree, so merely
permitting additional assumptions changes no rule instance. In a notation
that writes a finite permitted context at every node, freshen the local
parameters away from the additional context first; this gives the usual
explicit weakening rule with its freshness conditions.
For finite support, a finite tree has finitely many open leaves. Take
$\Gamma_0$ to be the set of their formulas. All leaves are still licensed
and all discharge and eigenparameter conditions remain the same.
\Done

For an infinite $\Gamma$, the notation $\Gamma\vdash B$ in this paragraph
means that a finite tree has its open leaves in $\Gamma$; it does not
require infinitely many unused assumptions to be written at every node.
Finite support is a syntactic fact about finitary proofs. It does not,
without a completeness theorem, establish semantic compactness.

\begin{example}[Reusing the inclusion proof]
The first displayed tree gives a derivation of $x\in B$ with open
assumptions $x\in A$ and $A\subseteq B$. Suppose separate derivations
establish $\Delta\vdash x\in A$ and $\Theta\vdash A\subseteq B$.
Grafting them at the corresponding leaves gives
\[
\Delta,\Theta\vdash x\in B.
\]
No new membership axiom is involved: the established derivations supply
the inputs of the existing proof. Likewise, grafting this result into
an inclusion proof with $B\subseteq C$ gives
$\Delta,\Theta,B\subseteq C\vdash x\in C$. This is composition of
proofs through their open assumptions, rather than substitution of
terms for free variables.
\end{example}

The occurrence bookkeeping is essential. An assumption introduced
inside a proof and later discharged cannot be replaced as though it
were an open dependency, and an inserted open assumption must not
accidentally acquire an old discharge label. For a calculus with quoted
sequents or reflection, each additional rule needs its own preservation
argument: replacing a derivation can change the context represented in
a quotation or the data tested by a reflection rule. The grafting theorem
above therefore does not assert cut admissibility for arbitrary encoded
reflection principles or for the earlier implementation merely because
they display the same sequent symbol.

\ifdefined\qwertyflag
\undef\qwertyflag
\end{document}
\fi
